\section{问题提出的背景}

\subsection{背景介绍}
近年来，以GPT-4\cite{gpt4}、Llama 3\cite{llama}为代表的大语言模型（LLM）在自然语言处理领域取得了里程碑式的突破。得益于Transformer\cite{transformer}架构的并行计算能力与海量预训练数据，LLM展现出了惊人的语言理解、生成与逻辑推理能力，被视为通向通用人工智能（AGI）的重要路径。然而，在实际应用中，LLM仍面临着一系列难以忽视的挑战：

具体而言，首先，LLM的知识来源受限于预训练语料，难以实时获取最新信息，且其基于概率的生成方式容易产生看似合理但事实错误的“幻觉”\cite{ref24}，在医疗、法律等高风险场景中会直接影响可靠性；其次，企业或个人的私有数据通常不包含在公开预训练数据中，通用模型难以直接适配垂直领域任务；最后，面对长尾与冷门知识时，即使模型参数规模很大，仍可能出现记忆不准确与概念混淆的问题。

为了解决上述问题，检索增强生成（Retrieval-Augmented Generation, RAG）范式应运而生。RAG通过在生成过程中引入外部非参数化知识库，使模型能够“查阅”最新、最准确的参考资料，从而显著提升了生成的可靠性\cite{ref1}。相较于微调（Fine-tuning），RAG具有成本低、知识更新快、可解释性强等优势，迅速成为LLM应用的主流架构。

然而，随着应用场景的深入，传统的RAG架构逐渐暴露出其局限性，特别是在处理复杂推理任务时：
主要体现在三个方面：其一，传统RAG多依赖单次语义检索，在科学研究或案件分析等任务中，答案往往分散于多文档的碎片证据，需要多步检索与多跳推理才能拼合证据链；其二，现实知识具有异构与多模态特征，文本、表格、知识图谱等信息源之间存在结构差异，单一文本检索难以有效利用图谱的结构逻辑或表格的数值关系，导致知识利用率受限；其三，尽管思维链（CoT）\cite{ref2}提示提升了模型的推理表现，但推理过程缺少外部证据的显式支撑，难以对关键结论进行可验证的溯源解释。

因此，如何构建一个能够融合多源异构知识，并具备复杂多跳推理能力的RAG系统，已成为当前学术界与工业界共同关注的焦点。

\textbf{核心概念界定：}为保证研究目标可落地、可评测，本文对核心概念作如下界定。
\textbf{（1）多源异构知识：}“多源”强调知识来源的多样性，包括公开语料（如百科/公开QA数据集）与私有/垂直领域数据（如企业文档、报告、PDF、表格等）；“异构”强调知识载体与结构的差异，既包括非结构化文本，也包括半结构化表格与结构化知识图谱，且同一事实可能以不同粒度、字段语义与版本时序的形式出现。本文所称“多源异构知识融合”并非简单拼接多路检索结果，而是要求在证据层面实现可定位、可对齐与可互证，使不同来源证据能够在同一推理过程中形成约束。
\textbf{（2）复杂推理：}本文将“复杂推理”限定为面向跨文档、多约束、可验证的推理过程，典型任务包括多跳问答与需要组合多个中间事实的检索增强推理。复杂推理不仅表现为生成更长的推理文本，更关键的是系统需要具备“证据选择—推理状态组织—验证与回退”的闭环控制能力，使结论能被证据链约束并可追溯。

\textbf{痛点与研究必要性：}在上述界定下，当前RAG系统的关键痛点可归纳为四类：其一，异构证据缺少统一的对齐与定位机制，导致跨载体证据难以稳定组合；其二，检索噪声与证据碎片化普遍存在，有限上下文预算下证据选择不稳定，容易把噪声放大为错误推理；其三，复杂推理链路缺少可验证机制，模型在证据不足或证据冲突时仍可能输出高置信结论，难以及时回退或拒答；其四，评测口径偏重最终答案指标，缺少对证据覆盖、检索/推理成本与证据支撑率等系统级过程指标的统一度量，导致迭代难以复现与归因。上述问题直接说明了本文研究的必要性：需要把多源异构融合与复杂推理视为一个统一系统问题，才能在效果、成本与可信性之间获得可解释且可持续迭代的折中。

\subsection{本研究的意义和目的}

本研究旨在探索多源异构知识融合的证据组织与协同推理机制，并在可运行的系统原型上进行验证。理论上，本研究尝试将基于连接主义的LLM与基于符号主义的知识图谱进行结合，探索神经符号人工智能在RAG架构下的可行范式；同时关注Graph of Thoughts（GoT）\cite{ref5}等非线性推理结构在异构知识库上的组织方式，推动推理从线性链式向更符合人类思维的网状结构演进；此外，本文将以“统一证据入口”为目标，研究文本、表格与图结构证据的表示、定位与对齐策略，在不依赖训练统一跨模态嵌入模型的前提下实现可落地的融合与推理。

在应用层面，本研究面向高复杂度的智能问答与决策支持需求，强调“可用证据支撑的可靠回答”。在科学研究辅助场景中，系统可帮助研究人员从文献、表格与结构化关系中快速发现潜在关联；在金融与法律分析中，系统可跨文档关联条款、记录与法规，形成可追溯的证据链；在企业知识管理中，系统可打破数据孤岛，将散落在文档与知识库中的信息整合为统一的检索与推理入口，提高组织知识利用效率。
